% Standalone problem document, generated from the adjacent README.md.
% Compile with XeLaTeX. Mathematical content is maintained in Markdown.
\documentclass[11pt,a4paper]{article}
\usepackage[fontsize=10.5pt]{scrextend}
\usepackage[margin=22mm,headheight=15pt,headsep=9mm,footskip=11mm]{geometry}
\usepackage{fontspec}
\setmainfont{texgyrepagella}[Extension=.otf,UprightFont=*-regular,BoldFont=*-bold,ItalicFont=*-italic,BoldItalicFont=*-bolditalic]
\setsansfont{texgyreheros}[Extension=.otf,UprightFont=*-regular,BoldFont=*-bold,ItalicFont=*-italic,BoldItalicFont=*-bolditalic]
\setmonofont{texgyrecursor}[Extension=.otf,UprightFont=*-regular,BoldFont=*-bold,ItalicFont=*-italic,BoldItalicFont=*-bolditalic,Scale=MatchLowercase]
\usepackage{amsmath,amssymb}
\usepackage{unicode-math}
\setmathfont{texgyrepagella-math.otf}
\usepackage{xcolor,microtype,enumitem,needspace,fancyhdr}
\usepackage{bookmark}
\definecolor{ink}{HTML}{17344B}
\definecolor{accent}{HTML}{197D80}
\definecolor{muted}{HTML}{596773}
\hypersetup{colorlinks=true,linkcolor=accent,urlcolor=accent,pdftitle={SP-06 - A
real-valued symbol on a Jordan curve and real Toeplitz
spectra},pdfauthor={Open Problems in Numerical Linear Algebra}}
\urlstyle{same}
\setlength{\parindent}{0pt}
\setlength{\parskip}{5pt plus 1pt minus 1pt}
\linespread{1.04}
\setlength{\emergencystretch}{3em}
\widowpenalty=10000
\clubpenalty=10000
\displaywidowpenalty=10000
\allowdisplaybreaks[1]
\setlist{nosep,leftmargin=1.5em}
\providecommand{\tightlist}{\setlength{\itemsep}{0pt}\setlength{\parskip}{0pt}}
\providecommand{\pandocbounded}[1]{#1}
\pagestyle{fancy}
\fancyhf{}
\fancyhead[L]{\sffamily\footnotesize\color{muted}OPEN PROBLEMS IN NUMERICAL LINEAR ALGEBRA}
\fancyhead[R]{\sffamily\bfseries\footnotesize\color{accent}SP-06}
\fancyfoot[L]{\parbox[b]{0.9\textwidth}{\sffamily\fontsize{8}{10}\selectfont\color{muted}Editorial ratings; a bounded search cannot certify that a problem remains unsolved.\\Literature check: 2026-09-13}}
\fancyfoot[R]{\sffamily\footnotesize\color{muted}\thepage}
\renewcommand{\headrulewidth}{0.3pt}
\renewcommand{\headrule}{\hbox to\headwidth{\color{accent}\leaders\hrule height \headrulewidth\hfill}}
\makeatletter
\renewcommand{\section}{\Needspace{5\baselineskip}\@startsection{section}{1}{0pt}{12pt plus 2pt minus 2pt}{4pt}{\sffamily\bfseries\large\color{ink}}}
\renewcommand{\subsection}{\Needspace{4\baselineskip}\@startsection{subsection}{2}{0pt}{9pt plus 1pt minus 1pt}{3pt}{\sffamily\bfseries\normalsize\color{ink}}}
\makeatother
\setcounter{secnumdepth}{0}
\begin{document}
{\sffamily\small\color{muted}Eigenvalues and inverse problems\par}
\vspace{3pt}
{\raggedright\hyphenpenalty=10000\exhyphenpenalty=10000\sffamily\bfseries\fontsize{21}{24}\selectfont\color{ink}A
real-valued symbol on a Jordan curve and real Toeplitz spectra\par}
\vspace{7pt}
{\sffamily\small\textbf{Difficulty:} challenging\quad\textbf{Importance:} interesting
to the community\par}
{\sffamily\small\bfseries\color{accent}Status: Lean verified\par}
\vspace{4pt}
{\color{accent}\hrule height 0.8pt}
\vspace{6pt}
\textbf{Topic:} Spectra of finite banded Toeplitz matrices.

\textbf{Rating rationale:} Challenging reflects recovering a spectral
implication after a gap in its original analytic proof; community impact
is a structural criterion for real spectra of banded Toeplitz sections.

\subsection{Resolution --- 2026-09-11}\label{resolution-2026-09-11}

\textbf{Negative resolution by Matthew J. Colbrook}, Department of
Applied Mathematics and Theoretical Physics, University of Cambridge.
\href{https://github.com/ajt60gaibb/OpenProblemsInNLA/blob/main/eigenvalues-and-inverse-problems/SP-06/solution.md}{Complete
manuscript} ·
\href{https://github.com/ajt60gaibb/OpenProblemsInNLA/blob/main/eigenvalues-and-inverse-problems/SP-06/solution.pdf}{PDF}
·
\href{https://github.com/ajt60gaibb/OpenProblemsInNLA/blob/main/eigenvalues-and-inverse-problems/SP-06/solution.tex}{LaTeX}.
\textbf{Theorem 1 and equations (1)--(8).}

The integer-coefficient Laurent polynomial in Theorem 1 is real on a
rigorously constructed star-shaped Jordan curve enclosing zero, while
its \(2\times2\) Toeplitz section has eigenvalues \(-128\pm8i\). The
proof checks continuity, injectivity and reality on the entire curve. It
refutes the conjectured implication for every finite section; it does
not refute the distinct limiting-spectrum statement.

The complete argument received an independent Codex-agent \textbf{PASS}
on 11 September 2026. The
\href{https://github.com/ajt60gaibb/OpenProblemsInNLA/blob/main/references/colbrook-additional-2026-09-11/verification/reviews/SP-06-review.md}{review
report} records the exact scope and a hash of the original reviewed
manuscript. The mathematical sections remain unchanged in the authored
version. Original ChatGPT generation is disclosed. That dated informal
review did not establish a formal proof certificate; the separate Lean
verification is recorded below. No external human peer review is
claimed.
\href{https://github.com/ajt60gaibb/OpenProblemsInNLA/blob/main/references/colbrook-additional-2026-09-11/README.md}{Submission
and verification record}.

The original statement, source references and prior audit notes are
retained; the former difficulty rating is historical.

\subsection{Lean proof and verification evidence ---
2026-09-13}\label{lean-proof-and-verification-evidence-2026-09-13}

\textbf{Formalization: George Stepaniants}, Department of Computing and
Mathematical Sciences, California Institute of Technology, Pasadena,
California, USA. Original mathematical proof credit remains Matthew J.
Colbrook, Department of Applied Mathematics and Theoretical Physics,
University of Cambridge.

The
\href{https://github.com/sgstepaniants/OpenProblemsInNLA/tree/3122d69460b0ed6dda3ea00dfaa899f93411ce2f/eigenvalues-and-inverse-problems/SP-06/lean}{immutable
complete proof} uses Lean 4.33.1, LeanCert
\textit{621a43d7cf21f87872392a01e874f2f1dbddc926}, and Mathlib
\textit{0df444a360eaa60ab8c11dca51a86af692955474}. The
\href{https://github.com/ajt60gaibb/OpenProblemsInNLA/blob/main/eigenvalues-and-inverse-problems/SP-06/lean/README.md}{project
guide},
\href{https://github.com/ajt60gaibb/OpenProblemsInNLA/blob/main/eigenvalues-and-inverse-problems/SP-06/lean/NUMERICAL_TARGETS.md}{exact
targets}, and
\href{https://github.com/ajt60gaibb/OpenProblemsInNLA/blob/main/eigenvalues-and-inverse-problems/SP-06/lean/SOURCE_MAP.md}{source
correspondence} identify all twenty exports and dependency pins.

\textit{NLA.SP06.witness\_counterexample} constructs the actual finite
Laurent polynomial, proves a continuous injective map of the entire
complex unit circle avoiding zero, proves the symbol real at every point
of its image, and exhibits a nonreal point of the actual complex
spectrum of its order-two Toeplitz section.
\textit{NLA.SP06.not\_targetImplication} negates the complete original
universal implication. The real Jordan curve is proved without assuming
an existence or continuity bridge. No restricted numerical sample
replaces the curve.

The repository ran
\href{https://github.com/sgstepaniants/OpenProblemsInNLA/actions/runs/34765629739/job/103745998196}{successful
Ubuntu verification} on the pinned proof revision. All twenty Comparator
statements, the permitted axiom closure and Lean's default kernel
passed, together with sandbox and rejection controls. Only
\textit{propext}, \textit{Classical.choice} and \textit{Quot.sound} are
permitted; the
\href{https://github.com/ajt60gaibb/OpenProblemsInNLA/blob/main/eigenvalues-and-inverse-problems/SP-06/lean/verification/local-2026-09-13/types-axioms.log}{per-export
axiom report} and
\href{https://github.com/ajt60gaibb/OpenProblemsInNLA/blob/main/eigenvalues-and-inverse-problems/SP-06/lean/verification/linux-2026-09-13/README.md}{raw
Linux evidence, independent audit and reproduction commands} are
retained. Two statement reviews preceded proof implementation, and two
independent final mathematical reviews passed. AI assistance and agent
review are disclosed; no human peer review or Tau Ceti endorsement is
claimed.

\subsection{Problem statement}\label{problem-statement}

Let \(r,s\ge1\) be integers and

\[b(z)=\sum_{k=-r}^{s}b_kz^k,
\qquad b_k\in\mathbb C,\qquad b_{-r}b_s\ne0.\]

For every \(n\ge1\), define \(T_n(b)=(b_{i-j})_{i,j=1}^n\), with
\(b_k=0\) outside \([-r,s]\). Suppose there is a Jordan curve
\(\gamma\subset\mathbb C\setminus\{0\}\) such that \(b(z)\in\mathbb R\)
for every \(z\in\gamma\). A Jordan curve means the image of a continuous
injective map from the unit circle.

Must

\[\mathop{\mathrm{spec}}\nolimits(T_n(b))\subset\mathbb R
\qquad\text{for every integer }n\ge1?\]

\subsection{Why it matters}\label{why-it-matters}

The conjecture would characterize a broad source of real spectra in
nonsymmetric structured eigenvalue problems through a scalar symbol
condition, uniformly over matrix size.

\newpage
\subsection{References}\label{references}

\begin{itemize}
\tightlist
\item
  B. Shapiro and F. Štampach,
  \href{https://arxiv.org/abs/1702.00741}{Non-Self-Adjoint Toeplitz
  Matrices Whose Principal Submatrices Have Real Spectrum},
  \emph{Constructive Approximation} 49 (2019), 191--226. Equation (2),
  Theorem 1(ii)--(iii), and Theorem 8; use arXiv v4, which includes the
  correction.
\item
  B. Shapiro and F. Štampach,
  \href{https://doi.org/10.1007/s00365-022-09614-0}{Correction to the
  same article}, 2023. The correction, pp.~27--28 of the combined arXiv
  PDF, withdraws the proof of (ii)\(\Rightarrow\)(iii) and explicitly
  proposes the implication as conjectural.
\item
  D. Giandinoto, \href{https://arxiv.org/abs/2411.16266}{On reality of
  eigenvalues of banded block Toeplitz matrices}, 2024, introduction and
  §2. The paper describes the corrected scalar implication as a
  conjecture before considering a block generalization.
\end{itemize}

\subsection{Earlier status check ---
2026-09-08}\label{earlier-status-check-2026-09-08}

On 2026-09-08, checked the corrected arXiv v4 (2022-12-29), its appended
erratum, and Giandinoto's November 2024 paper. Searches combined
``Shapiro'', ``Štampach'', ``Toeplitz'', ``reality'', ``erratum'',
``Jordan curve'', ``proof'', and 2026. No general resolution was
located. The original article's theorem label is misleading without its
correction: the finite-spectrum implication above remains the authors'
stronger proposed statement. Its weaker limiting-spectrum consequence
and particular symbol families are not separate catalog problems here.

\subsection{Audit update --- 2026-09-10}\label{audit-update-2026-09-10}

Rechecked the appended erratum in
\href{https://arxiv.org/pdf/1702.00741}{Shapiro--Štampach, version 4},
printed p.~27. It withdraws the relevant implication's proof and retains
it as a conjecture, so the original theorem wording must not be treated
as a solution. Jordan-curve/Toeplitz searches and
\href{https://arxiv.org/abs/2411.16266}{Giandinoto's later work} did not
identify a proof of this scalar sufficiency implication.
\end{document}
